\subsection{Catch-Me-If-You-Can: The Overshoot Problem and the Weak/Inflation Hierarchy --- arXiv:2207.00567}

\textbf{Authors:} Joseph P. Conlon, Filippo Revello

\paragraph{主な主張}
LVS の inflation 後に volume modulus が長い kination を進む間、初期には微小な radiation でも kinetic energy より遅く希釈されて tracker solution に追いつき、十分大きな inflation scale と weak scale の階層があれば minimum の overshoot を回避して安定真空へ到達できることを示した \cite{Conlon:2022catchme}。

\paragraph{新規性}
指数型ポテンシャルの tracker 自体は既知である一方、LVS 特有の長い field excursion を利用して seed radiation が tracker に成長する条件を定量化し、inflation/weak hierarchy を安定真空への宇宙論的到達可能性と結び付けた点にある。

\paragraph{位置付け}
string moduli の静的な真空安定化と宇宙論的な overshoot 問題を tracker dynamics で接続した具体例であり、一般の非指数型・多場モジュライポテンシャルで genuine tracker が成立するかを検討する際の基準となる。
